\documentclass[11pt,a4paper]{article}

% --- PAQUETES DE CONFIGURACIÓN ---
\usepackage[utf8]{inputenc}
\usepackage[T1]{fontenc}
\usepackage[spanish,es-tabla]{babel}
\usepackage{geometry}
\geometry{
    top=3.3cm,
    bottom=2.3cm,
    left=2.2cm,
    right=2.2cm,
    headheight=60pt,
    headsep=15pt,
    footskip=20pt
}

\usepackage{graphicx}
\usepackage{fancyhdr}
\usepackage{xcolor}
\usepackage{titlesec}
\usepackage{booktabs}
\usepackage{tabularx}
\usepackage{array}
\usepackage{amsmath}
\usepackage{enumitem}
\usepackage{microtype}
\usepackage{hyperref}

% --- PALETA INSTITUCIONAL FIEI / UNA PUNO ---
\definecolor{AzulUNA}{RGB}{10, 45, 90}
\definecolor{DoradoFIEI}{RGB}{180, 130, 20}
\definecolor{GrisOscuro}{RGB}{50, 50, 50}
\definecolor{RojoAlerta}{RGB}{150, 30, 30}

% --- CONFIGURACIÓN DE LISTAS ---
\setlist{topsep=2pt, itemsep=2pt, parsep=1pt}

% --- FORMATO DE ENLACES ---
\hypersetup{
    colorlinks=true,
    linkcolor=AzulUNA,
    urlcolor=AzulUNA,
    citecolor=AzulUNA
}

% --- FORMATO DE TÍTULOS ---
\titleformat{\section}{\color{AzulUNA}\normalfont\large\bfseries}{\thesection.}{0.7em}{}[\color{AzulUNA}\titlerule]
\titleformat{\subsection}{\color{AzulUNA}\normalfont\normalsize\bfseries}{\thesubsection.}{0.5em}{}
\titlespacing*{\section}{0pt}{7pt}{3pt}
\titlespacing*{\subsection}{0pt}{5pt}{2pt}

% --- ENCABEZADO Y PIE DE PÁGINA (fancyhdr) ---
\pagestyle{fancy}
\fancyhf{}

% Encabezado institucional: Logo FIEI (1.45cm) y Logo Equipo (2.05cm, 40% más grande)
\fancyhead[L]{%
    \includegraphics[height=1.45cm,keepaspectratio]{../../assets/img/logo_fiei.png}%
}
\fancyhead[C]{%
    \scriptsize \textbf{\textcolor{AzulUNA}{UNIVERSIDAD NACIONAL DEL ALTIPLANO DE PUNO}}\\[2pt]
    \tiny \textcolor{GrisOscuro}{FACULTAD DE INGENIERÍA ESTADÍSTICA E INFORMÁTICA}\\[2pt]
    \scriptsize \textbf{\textcolor{AzulUNA}{Gestión de Proyectos TICS (EST315) \textbar\ Team Sobrevivientes}}\\[2pt]
    \footnotesize \textbf{\textcolor{DoradoFIEI}{ACTA DE EVALUACIÓN DE CARTERA Y DESESTIMACIÓN TÉCNICA}}%
}
\fancyhead[R]{%
    \includegraphics[height=2.05cm,keepaspectratio]{../../assets/img/logo_equipo.png}%
}

\renewcommand{\headrulewidth}{0pt}
\renewcommand{\headrule}{}

\fancyfoot[L]{\scriptsize \textcolor{GrisOscuro}{Resolución de Cartera: Cancelación de Horus Sentinel}}
\fancyfoot[C]{\small \thepage}
\fancyfoot[R]{\scriptsize \textcolor{GrisOscuro}{Puno, Perú -- Semestre 2026-I}}
\renewcommand{\footrulewidth}{0.4pt}

\begin{document}

% =========================================================================
% PÁGINA 1: ANTECEDENTES Y SUSTENTO TÉCNICO DE LA DESESTIMACIÓN
% =========================================================================

\begin{center}
    {\Large \textbf{\textcolor{AzulUNA}{ACTA DE EVALUACIÓN TÉCNICA Y DESESTIMACIÓN DE PROYECTO}}}\\[3pt]
    {\normalsize \textbf{Resolución de Cancelación de Horus Sentinel y Priorización de CasScore}}\\[4pt]
    \footnotesize \textbf{Fecha de Sesión:} 08 de septiembre de 2026 \quad \textbar\quad \textbf{Estado:} Aprobado por Unanimidad
\end{center}

\vspace{0.1cm}

\noindent\begin{tabularx}{\textwidth}{@{}llXX@{}}
\toprule
\textbf{Curso:} EST315 -- Gestión de Proyectos TICS & \textbf{Docente:} Ing. Renzo Apaza Cutipa \\
\textbf{Equipo:} Team Sobrevivientes & \textbf{Decisión:} Cierre de Propuesta 2 / Ratificación de Propuesta 1 \\
\bottomrule
\end{tabularx}

\vspace{0.2cm}

\section{Antecedentes y Contexto de la Decisión}
Durante la etapa de exploración de la Unidad I, el equipo formuló dos propuestas tecnológicas: \textbf{CASA SCORE} (sistema de recomendación multicriterio para alquileres en Puno) y \textbf{Horus Sentinel} (plataforma geoespacial de mapeo predictivo de delincuencia). Tras la sustentación preliminar y la evaluación de la triple restricción (alcance, tiempo y costo), se convocó a sesión extraordinaria para evaluar la factibilidad técnica y operativa de continuar con ambas iniciativas o concentrar esfuerzos en una sola línea de desarrollo.

\section{Fundamentos Técnicos para la Desestimación de Horus Sentinel}
La propuesta \textit{Horus Sentinel} planteaba un alto valor conceptual enfocado en seguridad ciudadana; no obstante, el análisis de viabilidad técnica arrojó factores críticos e insalvables para el horizonte temporal del semestre:

\begin{enumerate}[leftmargin=*]
    \item \textbf{Inviabilidad en el Acceso y Disponibilidad de Datos Primarios:}
    El sistema dependía estrictamente de microdatos criminales georreferenciados (fecha, hora, coordenadas y tipología penal). Dichos registros son de acceso reservado por la Policía Nacional del Perú y el Ministerio del Interior. Tramitar convenios interinstitucionales de interoperabilidad o liberación de datos abiertos excede los plazos lectivos, generando un bloqueo insalvable para entrenar el modelo.

    \item \textbf{Excesiva Complejidad Algorítmica Espaciotemporal:}
    El modelado probabilístico requería articular tres técnicas de alta dificultad matemática:
    \begin{itemize}
        \item Estimación de Densidad Kernel (KDE) para el mapa de calor basal.
        \item Procesos Puntuales de Poisson No Homogéneos espaciotemporales para estimar réplicas delictivas.
        \item Algoritmo $A^*$ modificado sobre grafos viales con función biobjetivo: $\text{Costo} = \alpha \cdot \text{Distancia} + \beta \cdot \text{Riesgo}$.
    \end{itemize}
    Calibrar, validar y optimizar estos modelos sin distorsión demanda un plazo estimado no menor a 5 meses, incompatible con el ciclo de 6 semanas del curso.

    \item \textbf{Riesgo Crítico de Sesgo Policial (Efecto PredPol):}
    Sin fuentes diversificadas y auditadas, el algoritmo tiende a sobrevigilar cuadrantes donde históricamente se reportan más denuncias (no necesariamente donde ocurren más delitos), lo que compromete la ética algorítmica exigida en proyectos de ingeniería.

    \item \textbf{Costos de Infraestructura y Servidores Dedicados:}
    El procesamiento espacial continuo en tiempo real sobre grafos y mapas raster requiere recursos de cómputo en la nube que sobrepasan las capas gratuitas, contraviniendo la política de costo cero del equipo.
\end{enumerate}

\section{Matriz de Evaluación Crítica de Alternativas}
\vspace{0.1cm}
\noindent\begin{tabularx}{\textwidth}{p{4.2cm} X X}
\toprule
\textbf{Criterio de Viabilidad} & \textbf{Horus Sentinel (Desestimado)} & \textbf{CASA SCORE (Ratificado)} \\
\midrule
\textbf{Dependencia de Terceros} & \textbf{Crítica:} Depende de permisos de la PNP y Serenazgo. & \textbf{Nula:} Levantamiento directo en campo estudiantil. \\
\addlinespace[2pt]
\textbf{Curva de Desarrollo} & \textbf{Inviable:} Supera el semestre académico. & \textbf{Viable:} 5 sprints semanales para MVP. \\
\addlinespace[2pt]
\textbf{Riesgo de Fracaso en MVP} & \textbf{Muy Alto:} No se garantiza producto funcional en sem. 6. & \textbf{Bajo:} Alcance y arquitectura modular controlada. \\
\bottomrule
\end{tabularx}

\clearpage

% =========================================================================
% PÁGINA 2: RESOLUCIÓN FORMAL, REASIGNACIÓN Y SUSCRIPCIÓN
% =========================================================================

\section{Sustento de Viabilidad y Ratificación de CASA SCORE (AlquilaPuno)}
Frente al descarte de Horus Sentinel, el proyecto \textbf{CASA SCORE} demuestra las condiciones óptimas para su ejecución integral:
\begin{itemize}[leftmargin=*]
    \item \textbf{Impacto Social Universitario Inmediato:} Resuelve la asimetría y el riesgo de fraude habitacional que padecen los estudiantes que llegan a la UNA Puno.
    \item \textbf{Modelo Matemático Determinístico y Controlable:} La normalización lineal min-max en escala común $[0, 100]$ y la ponderación multicriterio garantizan rigor analítico verificable sin dependencia de bibliotecas externas complejas.
    \item \textbf{Stack Técnico Gratuito y Validado:} Operación 100\% en la nube mediante Next.js, Supabase (PostgreSQL), Leaflet y OpenStreetMap, permitiendo entregar un prototipo móvil/web en tiempo récord.
\end{itemize}

\section{Resoluciones y Acuerdos Formales del Equipo}
Estando presentes todos los integrantes de \textbf{Team Sobrevivientes}, se adoptan los siguientes acuerdos vinculantes:

\begin{enumerate}[leftmargin=*]
    \item \textbf{Primero (Cancelación Definitiva):} Declarar formalmente cancelado y archivado el proyecto \textit{Horus Sentinel}, congelando su ficha técnica únicamente como evidencia documental del análisis inicial de cartera de la Unidad I.
    \item \textbf{Segundo (Adopción de Proyecto Único):} Ratificar a \textbf{CASA SCORE} (bajo el nombre aplicativo de \textit{AlquilaPuno}) como el único proyecto oficial que desarrollará el equipo durante el ciclo académico 2026-I en el curso EST315.
    \item \textbf{Tercero (Reasignación de Cargas de Trabajo):} Reasignar el 100\% de las horas de trabajo e investigación técnica hacia el diseño de la base de datos relacional en Supabase, el algoritmo de ranking y la arquitectura del MVP.
    \item \textbf{Cuarto (Control en Repositorio):} La Gestora de Comunicación registrará este acuerdo en el repositorio GitHub, cerrando las ramas experimentales de Horus y procediendo a la apertura de la línea base del proyecto oficial.
\end{enumerate}

\vspace{0.4cm}

\section{Suscripción y Conformidad del Equipo}
\noindent En señal de pleno acuerdo con los fundamentos y resoluciones de la presente acta, los miembros del equipo firman a continuación:

\vspace{1.8cm}

% Bloque de firmas blindado con parbox centradas y espacio superior suficiente
\begin{center}
\begin{tabular}{c c c}
    \parbox{6.5cm}{\centering \rule{6.0cm}{0.6pt}\\[5pt] \textbf{Clyde Neil Paricahua Pari}\\[2pt] {\footnotesize Project Manager}} 
    & \hspace{0.8cm} &
    \parbox{6.5cm}{\centering \rule{6.0cm}{0.6pt}\\[5pt] \textbf{Ruth Karina Apaza Solis}\\[2pt] {\footnotesize Analista de Planificación}} 
    \\[2.2cm]
    \parbox{6.5cm}{\centering \rule{6.0cm}{0.6pt}\\[5pt] \textbf{Eliseo Tarqui Ajahuana}\\[2pt] {\footnotesize Especialista Técnico / Arquitectura}} 
    & \hspace{0.8cm} &
    \parbox{6.5cm}{\centering \rule{6.0cm}{0.6pt}\\[5pt] \textbf{Viviana Mary Quisbert Quispe}\\[2pt] {\footnotesize Gestora de Comunicación y Riesgos}}
\end{tabular}
\end{center}

\end{document}